% Generated draft: requires mathematical and source review.
\documentclass[11pt]{article}
\usepackage[margin=25mm]{geometry}
\usepackage{fontspec}
\setmainfont{Latin Modern Roman}
\usepackage{amsmath,amssymb,mathtools,unicode-math}
\setmathfont{Latin Modern Math}
\usepackage{enumitem,needspace,xurl,hyperref}
\setlength{\parindent}{0pt}
\setlength{\parskip}{5pt}
\setlength{\emergencystretch}{3em}
\newcommand{\N}{\mathbb{N}}
\newcommand{\Z}{\mathbb{Z}}
\newcommand{\Q}{\mathbb{Q}}
\newcommand{\R}{\mathbb{R}}
\newcommand{\C}{\mathbb{C}}
\newcommand{\F}{\mathbb{F}}
\newcommand{\A}{\mathbb{A}}
\newcommand{\PP}{\mathbb{P}}
\newcommand{\E}{\mathbb{E}}
\newcommand{\eps}{\varepsilon}
\newcommand{\dd}{\,\mathrm{d}}
\begin{document}
\section*{UnsolvedMath ID: 6200041}
\textbf{Source problem code:} catalogue entry
\subsection{Definitions and mathematical statement}
Let \(f:(X,d)\to(Y,\rho)\) be a homeomorphism. It is quasisymmetric if there is a homeomorphism \(\eta:[0,\infty)\to[0,\infty)\) such that for every triple of distinct points \(x,a,b\in X\),
\[
\frac{\rho(f(x),f(a))}{\rho(f(x),f(b))}
\le\eta\left(\frac{d(x,a)}{d(x,b)}\right).
\]
A common pointwise metric definition of quasiconformality bounds infinitesimal linear distortion. With
\[
L_f(x,r)=\sup_{d(x,a)\le r}\rho(f(x),f(a)),\qquad
l_f(x,r)=\inf_{d(x,b)\ge r}\rho(f(x),f(b)),
\]
for sufficiently small admissible radii, it requires
\[
H_f(x)=\limsup_{r\downarrow0}\frac{L_f(x,r)}{l_f(x,r)}\le K<\infty
\]
at every point. Other metric, almost-everywhere, analytic, and modulus definitions have different hypotheses; the source does not select one.

Kleiner asks for the generality of the implication quasiconformal \(\Rightarrow\) quasisymmetric. A precise classification would identify metric-space classes and the chosen quasiconformal notion for which every such homeomorphism satisfies the displayed three-point inequality, together with any requested dependence of \(\eta\) on \(K\) and geometric data. A local version asks for such control on neighborhoods instead of the entire space; these conclusions are distinct. Neither a global implication for arbitrary Euclidean domains nor a uniform distortion function is silently inferred from a local result.

The supplied research summary reports Euclidean local implications in the appropriate setting and sufficient hypotheses from Heinonen--Koskela theory involving Loewner and measure/doubling geometry. It gives no complete theorem hypotheses, so bare doubling alone is not promoted to a sufficient condition. The intended ambient measures, completeness, domain geometry, dimensional range, and necessity/classification target remain unspecified. The summary describes sufficient-condition progress, not a complete characterization.
\subsection{Short English statement}
Characterize the metric settings in which quasiconformal homeomorphisms must be quasisymmetric. The quasiconformal definition, local versus global conclusion, and geometric hypotheses must be fixed.
\subsection{Sources}
\begin{itemize}
\item Local UnsolvedMath registry, permanent ID \(6200041\), exact source code \texttt{catalogue entry}, ``Boundaries of Groups and Kleinian Groups --- Problem 41.'' Proposer Bruce Kleiner, 2005.
\item Kapovich, \emph{Problems on Boundaries of Groups and Kleinian Groups} (2005), Problem 41, recorded PDF page 12; \url{https://www.math.ucdavis.edu/~kapovich/EPR/problems.pdf}. Recorded access 2026-07-29; pdf-text-normalized extraction.
\item Supplied PARTIAL-PROGRESS research summary invokes Heinonen--Koskela theory, Loewner conditions, and doubling hypotheses without a complete publication citation or precise sufficient-condition theorem.
\end{itemize}
\end{document}
